
\section{Simple Nudging Example: 1D Heat Diffusion with a Single Observation}

\subsection*{Physical Model}

Consider the temperature evolution in a 1D rod of length $L$ governed by the heat equation:
\[
\frac{\partial T(x,t)}{\partial t} = \kappa \frac{\partial^2 T(x,t)}{\partial x^2}
\]
with homogeneous boundary conditions:
\[
T(0,t) = T(L,t) = 0, \quad T(x,0) = T_0(x)
\]

\subsection*{Observation Setup}

Assume we have one temperature sensor located at position $x = x_m$ that records the true temperature $T_{\text{obs}}(t) \approx T(x_m,t)$ over time.

\subsection*{Nudging-Augmented Model}

To incorporate the observation, we modify the equation by adding a nudging term:
\[
\frac{\partial T(x,t)}{\partial t} = \kappa \frac{\partial^2 T(x,t)}{\partial x^2} + \alpha (T_{\text{obs}}(t) - T(x,t)) \delta(x - x_m)
\]
where:
\begin{itemize}
  \item $\alpha$ is the nudging gain,
  \item $\delta(x - x_m)$ is a Dirac delta function centered at the sensor location.
\end{itemize}

\subsection*{Purpose and Interpretation}

This example demonstrates:
\begin{itemize}
  \item How localized observations influence a dynamical system,
  \item The impact of the nudging gain $\alpha$,
  \item How nudging helps reconcile simulated and observed data,
  \item That even a single observation can correct the full state over time via diffusion.
\end{itemize}

This is a pedagogical test case ideal for teaching nudging principles.

